\documentclass[../../2_specs]{subfiles}

\begin{document}

\chapter{QR Code Standard}

\section{Front Matter}

\subsection{Foreword}

% This chapter is a chaotic compilation and reinterpretation of the QR Code standard---inspired by the formatting of ISO/IEC 18004:2015 and ISO/IEC 18004:2024 (PNS ISO/IEC 18004:2025), as previewed by the Bureau of Philippine Standards of the Department of Trade and Industry of the Republic of the Philippines (DTI-BPS) and various sublicensed vendors (thanks for the sneak peeks, y’all). While the author did not have access to the full official standard, as they are expensive and very inaccessible, I pieced together this version through publicly available research papers, books, online articles, Wikipedia rabbit holes, and the likes.

% I also added annotations throughout, partly for my own understanding and to help other folks who want to get into QR code encoding/decoding without being legally or financially overwhelmed.

% This isn’t an official standard; It’s not ISO-blessed. This is a stitched-together knowledge quilt. But if it helps you understand how QR codes work, or if you’re building your own encoder/decoder (like I am!), then mission accomplished.

This chapter is an independent compilation and reinterpretation of the QR Code standard. 
It is informed by publicly available resources such as research papers, technical books, developer documentation, and collaborative knowledge bases. 
While I did not have access to the complete ISO/IEC publications due to their cost and limited availability, this document draws upon secondary sources and open references to provide an accessible explanation of QR Code structure and encoding principles.  

Where appropriate, annotations and commentary have been added to clarify complex topics and to make the material approachable for readers who may wish to implement their own encoder or decoder, or who are simply curious about how QR codes function.  

This is not an official or authoritative standard. It should be read as a study guide and commentary rather than as a normative reference.

Happy scanning, bestie \texttt{:3}

\subsection{Disclaimer}

This chapter is an independent, non-official reinterpretation of the QR Code standard, intended for educational and reference purposes. 
It is not affiliated with, endorsed by, or a substitute for any publication by the International Organization for Standardization (ISO), the International Electrotechnical Commission (IEC), or any related standards body.  

All technical descriptions herein are based on information from publicly accessible and legally obtainable sources.  Care has been taken to avoid reproducing verbatim text from proprietary publications.  Commentary and interpretation are original contributions of the author, provided to improve accessibility and understanding.  

This work is offered under principles of fair use and fair dealing for scholarly and educational purposes. Readers requiring normative or compliance-related guidance should consult the official ISO/IEC 18004 standard or its national equivalents.

\subsection{Introduction}

According to \textcite{wikipedia-qr}, a QR Code (Quick Response Code) is a two-dimensional \href{https://en.wikipedia.org/wiki/Barcode\#Matrix_(2D)_codes}{matrix
barcode} invented in 1994 by \href{https://en.wikipedia.org/wiki/Masahiro_Hara}{Masahiro Hara} of the \href{https://en.wikipedia.org/wiki/Japan}{Japanese} company \href{https://en.wikipedia.org/wiki/Denso\#DENSO_Wave}{Denso Wave}, originally for tracking automobile parts. 
QR Codes are recognizable by their square modules arranged on a contrasting background, typically with position-detecting \href{https://en.wikipedia.org/wiki/Fiducial_markers}{fiducial markers} in three corners.  

Data encoded within a QR Code is interpreted by imaging devices such as cameras, and recovered using error---correction algorithms—most notably \href{https://en.wikipedia.org/wiki/Reed\%E2\%80\%93Solomon_error_correction}{Reed--Solomon error correction}---to ensure robustness against noise and distortion. The encoded information is embedded in both the horizontal and vertical dimensions of the symbol, enabling compact and efficient data storage and retrieval.

\section{Definition of Terms}

\subsection{Abbreviations}

\begin{tabbing}
  2D \hspace{3em} \= two-dimensional \parencite{mw-2d} \\
  QR code \> quick-response code \parencite{wikipedia-qr} \\
  API \> Application Programming Interface \parencite{wikipedia-api} \\
\end{tabbing}

\section{Symbol Description}

QR code is a 2D barcode with the following characteristics:

\begin{enumerate}
  \item Formats:
    \begin{enumerate}[label=\alph*.]
      \item QR code, the full version; and
      \item Micro QR code, a smaller version of the QR code standard for applications where symbol size is limited.
    \end{enumerate}
  \item Character Set\parencite{wikipedia-qr}:
    \begin{enumerate}[label=\alph*.]
      \item Numeric: \texttt{0}, \texttt{1}, \texttt{2}, \texttt{3}, \texttt{4}, \texttt{5}, \texttt{6}, \texttt{7}, \texttt{8}, \texttt{9}
      \item Alphanumeric: \texttt{0}--\texttt{9}, \texttt{A}--\texttt{Z} (upper-case only), space, \texttt{\$} (dollar sign), \texttt{\%} (percent symbol), \texttt{*} (asterisk), \texttt{+} (plus symbol), \texttt{-} (dash), \texttt{.} (period), \texttt{/} (slash), \texttt{:} (colon)
      \item Binary/byte: \nameref{app:ISO-IEC-8859-1} as described in \hyperref[app:ISO-IEC-8859-1]{Appendix~\ref*{app:ISO-IEC-8859-1}}.
      \item Kanji/kana: \nameref{app:Shift-JIS-X-0208} as described in \hyperref[app:Shift-JIS-X-0208]{Appendix~\ref*{app:Shift-JIS-X-0208}}.
    \end{enumerate}
\end{enumerate}

\section{Modes}

\section{Character Sets}

\subsection{ISO/IEC 8859-1}

According to \textcite{wikipedia-ISO-8859-1}, ISO/IEC~8859-1:1998, \emph{Information technology---\href{https://en.wikipedia.org/wiki/8-bit_computing}{8-bit} single-\href{https://en.wikipedia.org/wiki/Byte}{byte} coded graphic \href{https://en.wikipedia.org/wiki/Character_(computing)}{character} sets---Part 1: Latin alphabet No.~1}, is part of the \href{https://en.wikipedia.org/wiki/ISO/IEC_8859}{ISO/IEC 8859} series of \href{https://en.wikipedia.org/wiki/ASCII}{ASCII}-based standard \href{https://en.wikipedia.org/wiki/Character_encoding}{character encodings}, first edition published in 1987. ISO/IEC~8859-1 encodes what it refers to as ``Latin alphabet no.~1'', consisting of 191 \href{https://en.wikipedia.org/wiki/Character_(computing)}{characters} from the \href{https://en.wikipedia.org/wiki/Latin_script}{Latin script}.


\subsubsection{Decode (Code → Character)}

\begin{enumerate}
    \item Start with the hex code. Example: \texttt{0x2F}
    \item Remove the prefix \texttt{0x}.
    \item Split into two nibbles:
    \begin{enumerate}[label=\alph*.]
        \item High nibble = first hex digit (\texttt{2}).
        \item Low nibble = second hex digit (\texttt{F}).
    \end{enumerate}
    \item Convert the high nibble into the row by multiplying it by 16 (hex shift). Example: \texttt{2} → \texttt{20}
    \item The low nibble is the column. Example: \texttt{F}.
    \item Locate the character at row \texttt{20}, column \texttt{F} in the table.
\end{enumerate}

\section{Versions}

\subsection{Explanation}

According to Denso Wave Incorporated, QR Code versions determine the data capacity and structure of a code\cite{denso-wave-qr-versions}.

The following tables are sourced from the official QR Code documentation provided by Denso Wave.

The figures in the Numeral, Alphanumeric, Binary and Kanji columns indicate the maximum allowable number of respective characters, including those for the data bit number.

For example, when using the Version 1 QR Code with correction level L, the maximum allowable characters for data bit number, numerals, binaries, and Kanji are 152, 25, 17, and 10 respectively.

As explained by Denso Wave in the aforementioned documentation, to determine the version of the QR code to use, suppose the data to be input consists of 100-digit numerals. This can be achieved by following the steps described below\cite{denso-wave-qr-versions}:

\begin{enumerate}
  \item In the QR Code table, choose "Numeric" as the type of input data.
  \item Choose an error correction level from the options of L, M, Q and H. For example, the error correction level that will be used is M.
  \item Find a figure in the table, 100 or over and the closest to 100 that is at the intersection with a correction level M row. The number of the version row that contains this figure is the most appropriate version number. We can see at the tables below that the version closest to that is Version 3 with capacity for up to 101 numeric characters at error correction M. 
\end{enumerate}

\end{document}